%! Rational Functions and Zeros — Unit 1, Lesson 1.8 — Homework
\documentclass[10pt]{article}
\usepackage{precalculus-article}
\usepackage{precalculus-boxes}
\newcommand{\TallMath}[1]{$\displaystyle #1\rule[-1.4em]{0pt}{3.2em}$}

\begin{document}

\pageheader{Unit 1, Lesson 1.8}{Homework}
\namedateperiod

\vspace{0.12in}

\begin{enumerate}[itemsep=14pt]

\item For each rational function below, state (i) the domain restriction(s) and
      (ii) all real zeros. Show your work.

      \begin{enumerate}[label=(\alph*), itemsep=10pt]
        \item $\displaystyle f(x) = \dfrac{x^2 - 9}{x + 5}$

              \vspace{0.06in}
              Domain restriction(s): \blank{3.5cm} \qquad Zeros: \blank{4cm}

        \item $\displaystyle g(x) = \dfrac{x^2 - 2x - 8}{x^2 - 4}$

              \vspace{0.06in}
              Domain restriction(s): \blank{3.5cm} \qquad Zeros: \blank{4cm}

        \item $\displaystyle h(x) = \dfrac{3x + 6}{x^2 + 1}$

              \vspace{0.06in}
              Domain restriction(s): \blank{3.5cm} \qquad Zeros: \blank{4cm}
      \end{enumerate}

\item Solve $\displaystyle\frac{x^2 - 1}{x - 3} \geq 0$ using a sign chart.
      Write the solution in interval notation.

      \medskip
      Critical values: \blank{5cm}

      \medskip
      \begin{center}
      \renewcommand{\arraystretch}{1.6}
      \begin{tabular}{|c|c|c|c|c|}
      \hline
      \textbf{Interval} & $(-\infty,\,-1)$ & $(-1,\,1)$ & $(1,\,3)$ & $(3,\,\infty)$ \\
      \hline
      Test value & \blank{1.5cm} & \blank{1.5cm} & \blank{1.5cm} & \blank{1.5cm} \\
      \hline
      Sign of $r$ & \blank{1.5cm} & \blank{1.5cm} & \blank{1.5cm} & \blank{1.5cm} \\
      \hline
      \end{tabular}
      \end{center}

      \vspace{0.06in}
      Solution: \blank{8cm}

\item In your own words, explain why the zeros of the denominator of a rational function
      are \textbf{not} zeros of the function, but they \textbf{still matter} when solving
      a rational inequality.

      \vspace{0.06in}
      \writelines{3}

\item A rational function $r$ has zeros at $x = 1$ and $x = -3$, and the denominator
      is zero at $x = 2$.

      The sign chart below is partially completed.
      Fill in the remaining signs, then solve $r(x) < 0$.

      \medskip
      \begin{center}
      \renewcommand{\arraystretch}{1.6}
      \begin{tabular}{|c|c|c|c|c|}
      \hline
      \textbf{Interval} & $(-\infty,\,-3)$ & $(-3,\,1)$ & $(1,\,2)$ & $(2,\,\infty)$ \\
      \hline
      Sign of $r$ & $(+)$ & \blank{1.5cm} & \blank{1.5cm} & $(+)$ \\
      \hline
      \end{tabular}
      \end{center}

      \vspace{0.06in}
      Solution of $r(x) < 0$: \blank{6cm}

\item \textbf{AP-Style Multiple Choice.} Which of the following is a zero of
      $\displaystyle r(x) = \frac{2x^2 + x - 3}{x - 5}$?

      \begin{enumerate}[label=(\Alph*), itemsep=3pt]
        \item $x = -\tfrac{3}{2}$
        \item $x = 1$
        \item $x = 5$
        \item $x = \tfrac{3}{2}$
      \end{enumerate}

\end{enumerate}

\vspace{0.1in}

\begin{reflectionbox}
\textbf{Lesson Preview.} In Lesson 1.9, we will investigate what happens to the graph of a
rational function \textit{near} the values where the denominator is zero --- discovering
vertical asymptotes and holes. Today's work on domain restrictions is the foundation.
\end{reflectionbox}

\end{document}
